
%\needspace{4\baselineskip}
\color{uniblue}{\subsection{Блокировка при неисправности цепей напряжения 2 секции шин (БНН 2сш)}}
\color{black}

\begin{enumerate}

\item
Блокировка при неисправности цепей напряжения используется для блокировки органов защит, которые могут работать неправильно при нарушениях в цепях напряжения второй секции шин. В устройстве блокировка при неисправности цепей напряжения ТН второй секции шин представлена функцией «БНН 2сш».

\item
Симметричность вторичных напряжений, подводимых к устройству, определяется уровнем напряжения обратной последовательности присоединения. Для предотвращения ложного блокирования защит и функций автоматики при несимметричных КЗ в первичной сети контролируется уровень тока обратной последовательности. Данный алгоритм регулируется уставкой SGF2 позволяя выявлять все виды несимметричных повреждений в цепях ТН. Срабатывание ИО данного алгоритма активирует выходной сигнал «БНН 2сш: Пуск U2I2» по истечении времени нерегулируемой уставки Т2.

\item
Если максимальное из трех линейных напряжений ниже уставки с введенным контролем, задаваемый уставкой SGF3, минимальный из трех фазных токов выше уставки, отсутствуют сигналы срабатывания измерительного органа ОПО функций ДЗ либо общего пуска МТЗ, активируется выходной сигнал «БНН 2сш: Пуск minU3p» по истечении времени нерегулируемой уставки Т3.

\item
Формирование общего сигнала срабатывания БНН «БНН 2сш: Срабатывание» происходит при наличии одного из сигналов:
\begin{itemize}
    \item выходной сигнал «БНН 2сш: Пуск U2I2»;
    \item выходной сигнал «БНН 2сш: Пуск minU3p»;
    \item выходной сигнал «БНН 2сш: Автомат ТН откл»;
    \item входной дискретный сигнал «Неиспр. ТН 2сш».
\end{itemize}

\item
В устройстве реализован контроль состояния автоматического выключателя во вторичных цепях ТН 2 секции шин при помощи соответствующих блок-контактов. Сигнал с блок-контактов АВ ТН 2 секции шин заводится на дискретный вход устройства «Автомат ТН 2сш». Для предотвращения кратковременных несимметричных режимов, возникающих при одновременном замыкании силовых контактов автомата, предусмотрена задержка на снятие блокировки ступеней защиты при включении автомата ТН длительностью 120 мс. Наличие аварийного или принудительного отключения АВ ТН формируется выходной сигнал «БНН 2сш: Автомат ТН откл».
\item
При наличии активного дискретного сигнала «Неиспр. ТН 2сш» и введенной функции в работу формируется сигнал общего срабатывания функции «БНН 2сш: Срабатывание».
\item
После появления сигнала «БНН 2сш: Срабатывание» запускается выдержка времени срабатывания сигнализации, задаваемая уставкой Т1, по истечении которой сформируется сигнализация о неисправности в цепях напряжения сигнал «БНН 2сш: Неисправность ЦН».

\item
Ввод функции в работу на этапе параметрирования устройства осуществляется программным переключателем SGF1 «Ввод\_функции». Информация о введенной ступени отображается выходным сигналом «БНН 2сш: Ввод» на выходе функции.
\item
Функция может быть выведена из работы оперативно путем активации сигнала «ОВ БНН 2сш», что характеризуется активным состоянием сигнала «БНН 2сш: Оперативный вывод».

\item
Алгоритм работы функции выполнен в соответствии с рисунком \ref{bnnsh2:img1}. В таблице \ref{bnnsh2:tbl1} приведены параметры, необходимые для настройки функции.

\vspace{3mm}
\begin{figure}[H]
\centering
\includegraphics[width=1\textwidth,height=1\textheight,keepaspectratio]{\fbpath/406. БНН 2сш/_latex/img/bnnsh2.pdf}
\captionof{figure}{Функциональная схема логики БНН 2сш}\label{bnnsh2:img1}
\end{figure}

\end{enumerate}
